\documentclass[10pt,a4paper]{amsart}
\usepackage{amsmath,amssymb,amsthm,graphicx}
% Retained preamble for the Stacks Project chapter snapshot in this directory.
% The chapter's own preamble.tex is a copy of the shared Stacks Project preamble; the block
% below is sliced MECHANICALLY, line for line, from it. Nothing was added, removed or
% redefined; the wrapper class, the xr-hyper/hyperref wiring and the externaldocument list
% are omitted because the assembler supplies the document class.
%
% lines taken from preamble.tex: 3, 5, 9, 11, 12, 156, 157, 158, ... (51 lines total)
%
% package coverage restored from the chapter's own preamble (2026-10-13): xy, xy-2cell

\usepackage[all]{xy}\usepackage{verbatim}
\xyoption{2cell}
\UseAllTwocells

\newenvironment{reference}{\comment}{\endcomment}
\newenvironment{slogan}{\comment}{\endcomment}
\newenvironment{history}{\comment}{\endcomment}
\theoremstyle{plain}
\newtheorem{theorem}[subsection]{Theorem}
\newtheorem{proposition}[subsection]{Proposition}
\newtheorem{lemma}[subsection]{Lemma}
\theoremstyle{definition}
\newtheorem{definition}[subsection]{Definition}
\newtheorem{example}[subsection]{Example}
\newtheorem{exercise}[subsection]{Exercise}
\newtheorem{situation}[subsection]{Situation}
\theoremstyle{remark}
\newtheorem{remark}[subsection]{Remark}
\newtheorem{remarks}[subsection]{Remarks}
\numberwithin{equation}{subsection}
\def\lim{\mathop{\mathrm{lim}}\nolimits}
\def\colim{\mathop{\mathrm{colim}}\nolimits}
\def\Spec{\mathop{\mathrm{Spec}}}
\def\Hom{\mathop{\mathrm{Hom}}\nolimits}
\def\Ext{\mathop{\mathrm{Ext}}\nolimits}
\def\SheafHom{\mathop{\mathcal{H}\!\mathit{om}}\nolimits}
\def\SheafExt{\mathop{\mathcal{E}\!\mathit{xt}}\nolimits}
\def\Sch{\mathit{Sch}}
\def\Mor{\mathop{\mathrm{Mor}}\nolimits}
\def\Ob{\mathop{\mathrm{Ob}}\nolimits}
\def\Sh{\mathop{\mathit{Sh}}\nolimits}
\def\NL{\mathop{N\!L}\nolimits}
\def\CH{\mathop{\mathrm{CH}}\nolimits}
\def\proetale{{pro\text{-}\acute{e}tale}}
\def\etale{{\acute{e}tale}}
\def\QCoh{\mathit{QCoh}}
\def\Ker{\mathop{\mathrm{Ker}}}
\def\Im{\mathop{\mathrm{Im}}}
\def\Coker{\mathop{\mathrm{Coker}}}
\def\Coim{\mathop{\mathrm{Coim}}}
\DeclareMathSymbol{\boxtimes}{\mathbin}{AMSa}{"02}
\def\QCohstack{\mathcal{QC}\!\mathit{oh}}
\def\Cohstack{\mathcal{C}\!\mathit{oh}}
\def\Spacesstack{\mathcal{S}\!\mathit{paces}}
\def\Quotfunctor{\mathrm{Quot}}
\def\Hilbfunctor{\mathrm{Hilb}}
\def\Curvesstack{\mathcal{C}\!\mathit{urves}}
\def\Polarizedstack{\mathcal{P}\!\mathit{olarized}}
\def\Complexesstack{\mathcal{C}\!\mathit{omplexes}}
\def\Pic{\mathop{\mathrm{Pic}}\nolimits}
\def\Picardstack{\mathcal{P}\!\mathit{ic}}
\def\Picardfunctor{\mathrm{Pic}}
\def\Deformationcategory{\mathcal{D}\!\mathit{ef}}
\usepackage{cleveref}
\crefname{originallemma}{original lemma}{original lemmas}
\Crefname{originallemma}{Original lemma}{Original lemmas}
\makeatletter
\newlabel{derived-lemma-prepare-adjoint}{{0CQR}{source}{}{sourcepointer-1}{}}
\newlabel{derived-lemma-prepare-adjoint@cref}{{[originallemma][0][]0CQR}{[1][0][]source}{}{}{}}
\makeatother
\crefname{originallemma}{original lemma}{original lemmas}
\Crefname{originallemma}{Original lemma}{Original lemmas}
\makeatletter
\newlabel{derived-lemma-right-adjoint}{{0CQS}{source}{}{sourcepointer-2}{}}
\newlabel{derived-lemma-right-adjoint@cref}{{[originallemma][0][]0CQS}{[1][0][]source}{}{}{}}
\makeatother
\crefname{originallemma}{original lemma}{original lemmas}
\Crefname{originallemma}{Original lemma}{Original lemmas}
\makeatletter
\newlabel{lemma-affine-coherator}{{08H0}{source}{}{sourcepointer-3}{}}
\newlabel{lemma-affine-coherator@cref}{{[originallemma][0][]08H0}{[1][0][]source}{}{}{}}
\makeatother
\crefname{originallemma}{original lemma}{original lemmas}
\Crefname{originallemma}{Original lemma}{Original lemmas}
\makeatletter
\newlabel{lemma-exact-sequence-j-star}{{08GW}{source}{}{sourcepointer-4}{}}
\newlabel{lemma-exact-sequence-j-star@cref}{{[originallemma][0][]08GW}{[1][0][]source}{}{}{}}
\makeatother
\crefname{originallemma}{original lemma}{original lemmas}
\Crefname{originallemma}{Original lemma}{Original lemmas}
\makeatletter
\newlabel{lemma-induction-principle}{{08GP}{source}{}{sourcepointer-5}{}}
\newlabel{lemma-induction-principle@cref}{{[originallemma][0][]08GP}{[1][0][]source}{}{}{}}
\makeatother
\crefname{originallemma}{original lemma}{original lemmas}
\Crefname{originallemma}{Original lemma}{Original lemmas}
\makeatletter
\newlabel{lemma-quasi-coherence-direct-image}{{08FA}{source}{}{sourcepointer-6}{}}
\newlabel{lemma-quasi-coherence-direct-image@cref}{{[originallemma][0][]08FA}{[1][0][]source}{}{}{}}
\makeatother
\crefname{originallemma}{original lemma}{original lemmas}
\Crefname{originallemma}{Original lemma}{Original lemmas}
\makeatletter
\newlabel{lemma-quasi-coherence-pullback}{{08F4}{source}{}{sourcepointer-7}{}}
\newlabel{lemma-quasi-coherence-pullback@cref}{{[originallemma][0][]08F4}{[1][0][]source}{}{}{}}
\makeatother
\crefname{originallemma}{original lemma}{original lemmas}
\Crefname{originallemma}{Original lemma}{Original lemmas}
\makeatletter
\newlabel{section-coherator}{{08GX}{source}{}{sourcepointer-8}{}}
\newlabel{section-coherator@cref}{{[originallemma][0][]08GX}{[1][0][]source}{}{}{}}
\makeatother
\title{Stacks, Derived Categories of Spaces: better-coherator}
\author{The Stacks Project Authors}
\date{stacks-project:spaces-perfect.tex:local-snapshot-20261012; GNU FDL 1.2 or later}
\begin{document}
\maketitle
\noindent\textit{Editorial note.} Statement span is lines 4653-4659 of sources/stacks-spaces-perfect/pinned-original-source.tex; proof span is lines 4661-4730, the complete author proof (68 body lines).\par\medskip
\noindent\textit{Editorial note.} Navigation stubs: 8 cited-but-not-included labels are supplied as non-load-bearing pointers with their pinned tags where the snapshot tag index has them.\par\medskip
\noindent\textit{Editorial note.} extra\_packages is empty on purpose: the retained preamble already loads xy with options and a second declaration clashes in the assembler wrapper.\par\medskip
\noindent\textit{Editorial note.} Text is retained verbatim, including the author's own typos and placeholders.\par\medskip
\noindent\textit{Editorial note.} Queue index 30 of queue\_spaceperfect.json (pb=2661, stmt 4653-4659, proof 4661-4730).\par\medskip
\noindent\textit{Editorial note.} stacks-spaces-perfect wave sp01 (2026-10-14): fresh adjudication accept per lane-s3/sp01\_judgments; pool per Chief 12:4x unfreeze (reserve-pool workflow, six-check clean); spaces-perfect.tex 126 with-proof rows, 70 b107-occupied labels avoided, 56 free rows judged 32 accept / 24 reject.\par\medskip
\noindent\textit{Editorial note.} Cross-chapter same-label warning (NOT an occupancy conflict, Chief ruling 2026-10-13 evening): label 'lemma-better-coherator' also exists in stacks-perfect — a different proposition in a different chapter file.\par\medskip
\noindent\textit{Allowed standard prerequisites (editorial).}\begin{itemize}
\item Explicit prerequisite boundary. The cited inputs are prerequisites; the argument in the retained span is not replaced by any cited result.
\end{itemize}\medskip
\ifdefined\hypertarget\hypertarget{sourcepointer-1}{}\fi
\noindent\textit{Source reference.} Original lemma 0CQR: sources/stacks-spaces-perfect/pinned-original-source.tex, source lines 4730--4730. Cited by the retained statement or proof but not part of the collected spans; supplied only as a navigation pointer. The argument does not depend on its content.\par\medskip
\ifdefined\hypertarget\hypertarget{sourcepointer-2}{}\fi
\noindent\textit{Source reference.} Original lemma 0CQS: sources/stacks-spaces-perfect/pinned-original-source.tex, source lines 4730--4730. Cited by the retained statement or proof but not part of the collected spans; supplied only as a navigation pointer. The argument does not depend on its content.\par\medskip
\ifdefined\hypertarget\hypertarget{sourcepointer-3}{}\fi
\noindent\textit{Source reference.} Original lemma 08H0: sources/stacks-spaces-perfect/pinned-original-source.tex, source lines 2441--2454. Cited by the retained statement or proof but not part of the collected spans; supplied only as a navigation pointer. The argument does not depend on its content.\par\medskip
\ifdefined\hypertarget\hypertarget{sourcepointer-4}{}\fi
\noindent\textit{Source reference.} Original lemma 08GW: sources/stacks-spaces-perfect/pinned-original-source.tex, source lines 1889--1909. Cited by the retained statement or proof but not part of the collected spans; supplied only as a navigation pointer. The argument does not depend on its content.\par\medskip
\ifdefined\hypertarget\hypertarget{sourcepointer-5}{}\fi
\noindent\textit{Source reference.} Original lemma 08GP: sources/stacks-spaces-perfect/pinned-original-source.tex, source lines 1563--1577. Cited by the retained statement or proof but not part of the collected spans; supplied only as a navigation pointer. The argument does not depend on its content.\par\medskip
\ifdefined\hypertarget\hypertarget{sourcepointer-6}{}\fi
\noindent\textit{Source reference.} Original lemma 08FA: sources/stacks-spaces-perfect/pinned-original-source.tex, source lines 790--804. Cited by the retained statement or proof but not part of the collected spans; supplied only as a navigation pointer. The argument does not depend on its content.\par\medskip
\ifdefined\hypertarget\hypertarget{sourcepointer-7}{}\fi
\noindent\textit{Source reference.} Original lemma 08F4: sources/stacks-spaces-perfect/pinned-original-source.tex, source lines 620--625. Cited by the retained statement or proof but not part of the collected spans; supplied only as a navigation pointer. The argument does not depend on its content.\par\medskip
\ifdefined\hypertarget\hypertarget{sourcepointer-8}{}\fi
\noindent\textit{Source reference.} Original lemma 08GX: sources/stacks-spaces-perfect/pinned-original-source.tex, source lines 2300--2353. Cited by the retained statement or proof but not part of the collected spans; supplied only as a navigation pointer. The argument does not depend on its content.\par\medskip
\begin{lemma}
\label{lemma-better-coherator}
Let $S$ be a scheme.
Let $X$ be a quasi-compact and quasi-separated algebraic space over $S$.
The inclusion functor $D_\QCoh(\mathcal{O}_X) \to D(\mathcal{O}_X)$
has a right adjoint.
\end{lemma}

\begin{proof}[First proof]
We will use the induction principle in
Lemma \ref{lemma-induction-principle}
to prove this. If $D(\QCoh(\mathcal{O}_X)) \to D_\QCoh(\mathcal{O}_X)$
is an equivalence, then the lemma is true because the functor
$RQ_X$ of Section \ref{section-coherator} is a right adjoint to the functor
$D(\QCoh(\mathcal{O}_X)) \to D(\mathcal{O}_X)$.
In particular, our lemma is true for affine algebraic spaces, see
Lemma \ref{lemma-affine-coherator}.
Thus we see that it suffices to show: if $(U \subset X, f : V \to X)$
is an elementary distinguished square with $U$ quasi-compact
and $V$ affine and the lemma holds for $U$, $V$, and $U \times_X V$,
then the lemma holds for $X$.

\medskip\noindent
The adjoint exists if and only if for every object $K$ of
$D(\mathcal{O}_X)$ we can find a distinguished triangle
$$
E' \to E \to K \to E'[1]
$$
in $D(\mathcal{O}_X)$
such that $E'$ is in $D_\QCoh(\mathcal{O}_X)$ and such that
$\Hom(M, K) = 0$ for all $M$ in $D_\QCoh(\mathcal{O}_X)$. See
Derived Categories, Lemma \ref{derived-lemma-right-adjoint}.
Consider the distinguished triangle
$$
E \to Rj_{U, *}E|_U \oplus Rj_{V, *}E|_V \to
Rj_{U \times_X V, *}E|_{U \times_X V} \to E[1]
$$
in $D(\mathcal{O}_X)$ of Lemma \ref{lemma-exact-sequence-j-star}.
By Derived Categories, Lemma \ref{derived-lemma-prepare-adjoint}
it suffices to construct the desired distinguished triangles
for $Rj_{U, *}E|_U$, $Rj_{V, *}E|_V$, and
$Rj_{U \times_X V, *}E|_{U \times_X V}$. This reduces us to the statement
discussed in the next paragraph.

\medskip\noindent
Let $j : U \to X$ be an \'etale morphism corresponding with
$U$ quasi-compact and quasi-separated and the lemma is true for $U$.
Let $L$ be an object of $D(\mathcal{O}_U)$.
Then there exists a distinguished triangle
$$
E' \to Rj_*L \to K \to E'[1]
$$
in $D(\mathcal{O}_X)$
such that $E'$ is in $D_\QCoh(\mathcal{O}_X)$ and such that
$\Hom(M, K) = 0$ for all $M$ in $D_\QCoh(\mathcal{O}_X)$.
To see this we choose a distinguished triangle
$$
L' \to L \to Q \to L'[1]
$$
in $D(\mathcal{O}_U)$ such that $L'$ is in $D_\QCoh(\mathcal{O}_U)$
and such that $\Hom(N, Q) = 0$ for all $N$ in $D_\QCoh(\mathcal{O}_U)$.
This is possible because the statement in
Derived Categories, Lemma \ref{derived-lemma-right-adjoint}
is an if and only if.
We obtain a distinguished triangle
$$
Rj_*L' \to Rj_*L \to Rj_*Q \to Rj_*L'[1]
$$
in $D(\mathcal{O}_X)$. Observe that $Rj_*L'$ is in $D_\QCoh(\mathcal{O}_X)$
by Lemma \ref{lemma-quasi-coherence-direct-image}.
On the other hand, if $M$ in $D_\QCoh(\mathcal{O}_X)$, then
$$
\Hom(M, Rj_*Q) = \Hom(Lj^*M, Q) = 0
$$
because $Lj^*M$ is in $D_\QCoh(\mathcal{O}_U)$ by
Lemma \ref{lemma-quasi-coherence-pullback}.
This finishes the proof.
\end{proof}

\end{document}
